\documentclass[11pt]{amsart}

\usepackage{amssymb}
\usepackage{mathtools}
\usepackage[hidelinks,
  pdftitle={Order-zero operators and irreducibility of power-related matrix weights},
  pdfauthor={Knowledge Sense},
  pdfsubject={Matrix-valued orthogonal polynomials},
  pdfkeywords={matrix-valued orthogonal polynomials, irreducible weights, algebra D(W), operators of order zero}
  ]{hyperref}

\newtheorem{theorem}{Theorem}[section]
\newtheorem{proposition}[theorem]{Proposition}
\newtheorem{lemma}[theorem]{Lemma}
\newtheorem{corollary}[theorem]{Corollary}
\theoremstyle{definition}
\newtheorem{definition}[theorem]{Definition}
\theoremstyle{remark}
\newtheorem{remark}[theorem]{Remark}

\numberwithin{equation}{section}

\hyphenation{manuscript}
\setlength{\emergencystretch}{1.5em}

\newcommand{\C}{\mathbb{C}}
\newcommand{\R}{\mathbb{R}}
\newcommand{\N}{\mathbb{N}}
\newcommand{\Mat}{\operatorname{Mat}}
\newcommand{\diag}{\operatorname{diag}}
\newcommand{\cD}{\mathcal{D}}
\newcommand{\CR}{\mathcal{C}_{\R}}
\newcommand{\sL}{\mathsf{L}}
\newcommand{\sE}{\mathsf{E}}
\newcommand{\ip}[2]{\langle #1,#2\rangle}
\newcommand{\conj}[1]{\overline{#1}}

\begin{document}

\title[Power-related matrix weights]{Order-zero operators and irreducibility of power-related matrix weights}

\author{Knowledge Sense}
\address{Knowledge Sense}
\email{math@knowledgesense.jp}

\date{13 October 2026}

\subjclass[2020]{Primary 42C05; Secondary 33C45, 15A21}

\keywords{Matrix-valued orthogonal polynomials, irreducible weight matrices, algebra $\cD(W)$, operators of order zero}

\begin{abstract}
Weight matrices $W=T\diag(w_1,\dots,w_N)T^*$ of every size $N$, with
$T(x)=I+xA$ for an explicit nilpotent matrix $A$ depending on nonzero real
parameters, were introduced by Bono Parisi and Pacharoni for certain
classical scalar weights and studied by Bono Parisi for arbitrary scalar
weights. Bono Parisi conjectured that, when the scalar weights are related by
powers of $x$, the algebra $\cD(W)$ contains only scalar operators of order
zero, and he stated that this holds for $N\le 10$. The power relation can be
read in two ways, giving $w_j=w\,x^{s_j}$ with $(s_j)=(0,0,1,\dots,N-2)$ or
$(s_j)=(0,1,\dots,N-1)$. We prove the conjecture for both exponent lists and
every $N$. More precisely, for all nonzero real parameters and every
admissible scalar weight $w$ on an interval, bounded or unbounded, the
operators of order zero in $\cD(W)$ are the complex multiples of the
identity, and $W$ is irreducible. Since the interval may be unbounded, we
show directly that for weights $W=wP$, with $w$ a scalar weight and $P$ a
polynomial matrix, the existence of a non-scalar operator of order zero
implies the existence of a non-scalar $W$-symmetric constant matrix. As a
consequence, the Laguerre-type examples give, in every size, irreducible
weight matrices that are Darboux transformations of diagonal scalar weights
and whose orthogonal polynomials satisfy a second-order differential
equation.
\end{abstract}

\maketitle

\section{Introduction}

A weight matrix of size $N$ on an interval $\Delta=(x_0,x_1)$, with
$-\infty\le x_0<x_1\le\infty$, is an integrable function
$W\colon\R\to\Mat_N(\C)$ that is positive definite almost everywhere on
$\Delta$, vanishes outside $\Delta$, and has finite moments of all orders.
For a sequence $(Q_n)_{n\ge0}$ of matrix-valued orthogonal polynomials for
$W$, the algebra $\cD(W)$ consists of the differential operators
$D=\sum_{j=0}^m\partial^jF_j(x)$, acting on the right by
$P\cdot D=\sum_j\partial^j(P)F_j$, such that
$Q_n\cdot D=\Lambda_n(D)Q_n$ with $\Lambda_n(D)\in\Mat_N(\C)$ for every $n$
\cite{GT07}, \cite[Definition~2.2]{BP26}; it does not depend on the choice
of the sequence. A weight matrix is \emph{reducible} if there are a
nonsingular $M\in\Mat_N(\C)$ and weight matrices $W_1,W_2$ of positive sizes
with $W=M\diag(W_1,W_2)M^*$ on $\Delta$, and \emph{irreducible} otherwise
\cite[Definition~2.4]{BP26}. Tirao and Zurri\'an \cite{TZ18} showed that $W$
is reducible if and only if the real vector space
\[
\CR(W)=\{F\in\Mat_N(\C)\colon FW(x)=W(x)F^*\text{ for almost every }x\in\Delta\}
\]
of \emph{$W$-symmetric} matrices is larger than $\R I$, and that a reducible
weight has non-scalar operators of order zero in $\cD(W)$; the converse part
of \cite[Theorem~4.3]{TZ18} assumes a bounded interval.

Let $N\ge1$ and let $a_1,\dots,a_{N-1}$ be nonzero real numbers. With the
standard matrix units $E_{ij}$ put
\begin{equation}\label{eq:A}
A=\sum_{r=1}^{\lfloor N/2\rfloor}a_{2r-1}E_{2r-1,2r}
 +\sum_{r=1}^{\lfloor (N-1)/2\rfloor}a_{2r}E_{2r+1,2r},
\end{equation}
and $T(x)=e^{xA}$ for real $x$. Every nonzero entry of $A$ has an odd row index and an even column index, so
$A^2=0$, $T(x)=I+xA$ and $T(x)^{-1}=I-xA$. Weight matrices of the form
\begin{equation}\label{eq:BPweight}
W(x)=T(x)\diag\bigl(w_1(x),\dots,w_N(x)\bigr)T(x)^*
\end{equation}
were introduced by Bono Parisi and Pacharoni \cite[Definition~1.1]{BPP} for
classical scalar weights of the same type whose quotients $w_i/w_j$ are not
rational functions. Bono Parisi \cite{BP26} studied \eqref{eq:BPweight} for
arbitrary scalar weights $w_1,\dots,w_N$ on a common interval. He gave
explicit orthogonal polynomials for $W$, sufficient conditions for
bispectrality, and several irreducibility results. For $N=2$ the weight
reduces if and only if $w_1(x)=-a_1^2w_2(x)(b-x)(c-x)$ for some real $b<c$
with $\Delta\subseteq(b,c)$ \cite[Proposition~5.1]{BP26}; for arbitrary $N$ it is
irreducible when no quotient $w_i/w_j$ with $i\ne j$ is a rational function
\cite[Proposition~5.6]{BP26}, a statement attributed there to
\cite[Proposition~3.7]{BPP}, where it is proved for the Hermite-, Laguerre-
and Jacobi-type families of \cite[Definition~1.1]{BPP}. At the end of
\cite[Section~5]{BP26} the weights are related instead by powers of $x$,
$w_{i+1}(x)=w_1(x)x^{i-1}$, so that all these quotients are rational. It is
stated there that for $N\le10$ the algebra $\cD(W)$ contains only scalar
operators of order zero, and conjectured that this holds for arbitrary $N$.
Read with $i=1,\dots,N-1$, the displayed relation gives the exponents
$(0,0,1,\dots,N-2)$, whereas the $10\times10$ Laguerre-type example that
follows it, with $w_j(x)=e^{-x}x^{\alpha+j-1}$, has the exponents
$(0,1,\dots,9)$. We prove the conjecture for both exponent lists.

For $N\ge1$ let
\[
\sL_N=(0,0,1,2,\dots,N-2)\ \ (N\ge2),\quad \sL_1=(0),\quad
\sE_N=(0,1,2,\dots,N-1).
\]
A scalar weight is a weight matrix of size one. For a list
$s=(s_1,\dots,s_N)$ of nonnegative integers we call a scalar weight $w$
\emph{$s$-admissible} on $\Delta$ if $w(x)x^{s_j}$ is a scalar weight on
$\Delta$ for every $j$ (with $x^0=1$). We put
\begin{equation}\label{eq:Ws}
\begin{gathered}
D_s(x)=\diag(x^{s_1},\dots,x^{s_N}),\qquad P_s(x)=T(x)D_s(x)T(x)^*,\\
W_s(x)=w(x)P_s(x).
\end{gathered}
\end{equation}
Thus $W_s$ is the weight \eqref{eq:BPweight} with $w_j=w\,x^{s_j}$. Since
$s_1=0$ for both exponent lists, $w$ itself is a scalar weight. If $s$ contains an odd
exponent, which happens for $\sE_N$ with $N\ge2$ and for $\sL_N$ with
$N\ge3$, admissibility forces $x_0\ge0$; for $\sL_2=(0,0)$ and $N=1$ any
interval is possible. In all cases $W_s$ is a weight matrix on $\Delta$:
it is positive definite wherever all $w(x)x^{s_j}$ are positive, because
$T(x)$ is invertible, and its entries are $w$ times polynomials. Write
$\cD_0(W)$ for the set of operators of order zero in $\cD(W)$.

\begin{theorem}\label{thm:main}
Let $N\ge1$, let $a_1,\dots,a_{N-1}$ be nonzero real numbers, let
$s\in\{\sL_N,\sE_N\}$, and let $w$ be an $s$-admissible scalar weight on an
interval $\Delta$, bounded or not. Then:
\begin{enumerate}
\item[(a)] $\CR(W_s)=\R I$;
\item[(b)] $W_s$ is irreducible;
\item[(c)] $\cD_0(W_s)=\C I$, that is, every operator of order zero in
$\cD(W_s)$ is a complex multiple of the identity.
\end{enumerate}
\end{theorem}

Thus the conjecture at the end of \cite[Section~5]{BP26} holds for every
$N$ and for both exponent lists.

\begin{corollary}\label{cor:conjecture}
Let $N\ge1$, let $a_1,\dots,a_{N-1}$ be nonzero real numbers, and let
$w_1,\dots,w_N$ be scalar weights on a common interval with either $w_{i+1}(x)=w_1(x)x^{i-1}$ for $1\le i\le N-1$, or
$w_{i+1}(x)=w_1(x)x^{i}$ for $1\le i\le N-1$. Then the weight matrix
\eqref{eq:BPweight} is irreducible and $\cD(W)$ contains only scalar
operators of order zero.
\end{corollary}

Combined with the Laguerre-type results of \cite[Section~4.1]{BP26},
Theorem~\ref{thm:main} gives irreducible examples of every size with the
following properties.

\begin{corollary}\label{cor:laguerre}
Let $N\ge1$, $\alpha>-1$, $s\in\{\sL_N,\sE_N\}$, and let $a_1,\dots,a_{N-1}$
be nonzero real numbers. The weight matrix
\[
W(x)=T(x)\diag\bigl(e^{-x}x^{\alpha+s_1},\dots,e^{-x}x^{\alpha+s_N}\bigr)
T(x)^*,\qquad x\in(0,\infty),
\]
is irreducible and satisfies $\cD_0(W)=\C I$. The orthogonal polynomials of
$W$ given in \cite[Theorem~3.1]{BP26} are eigenfunctions of a second-order differential
operator, and $W$ is a Darboux transformation of
$\diag(e^{-x}x^{\alpha+s_1},\dots,e^{-x}x^{\alpha+s_N})$ in the sense of
\cite[Definition~2.6]{BP26}.
\end{corollary}

\begin{proof}
The function $w(x)=e^{-x}x^\alpha$ on $(0,\infty)$ is $s$-admissible because
$\alpha+s_j>-1$, so the first assertion is Theorem~\ref{thm:main}. The
second-order operator is the one constructed in \cite[Section~4.1]{BP26}
from \cite[Theorem~4.1]{BP26}, which applies to arbitrary Laguerre
parameters. The parameters $\alpha+s_j$ have integral differences, so
\cite[Theorem~4.3]{BP26} gives the Darboux transformation.
\end{proof}

For $N=10$ and $s=\sE_{10}$ this is the example at the end of
\cite[Section~5]{BP26}. Irreducible weights \eqref{eq:BPweight} of every size
were known from the criterion on non-rational quotients
\cite[Proposition~5.6]{BP26}; for the Hermite-, Laguerre- and Jacobi-type
families of \cite[Definition~1.1]{BPP} those weights are not Darboux
transformations of direct sums of classical scalar weights in the sense of
\cite[Definition~2.5]{BPP}, by \cite[Theorem~1.4]{BPP}; see also
\cite[Section~4]{BP26}. Irreducible Laguerre-type weights that are Darboux
transformations of diagonal scalar weights in the sense of
\cite[Definition~2.6]{BP26} were exhibited in \cite{BP26} for
$N=5$, with exponents $(0,0,1,1,2)$ \cite[(4.3) and Section~5]{BP26}, and in
the power-related cases stated for $N\le10$; Corollary~\ref{cor:laguerre}
provides them in every size.

\subsection*{Known and new parts}
The construction \eqref{eq:A}--\eqref{eq:BPweight} was introduced by Bono
Parisi and Pacharoni \cite{BPP} for the classical families of
\cite[Definition~1.1]{BPP} and extended to arbitrary scalar weights by Bono
Parisi \cite{BP26}. The power-related weights, the conjecture and the
statement for $N\le10$ are due to Bono Parisi \cite{BP26}; the case $N=2$ of
the irreducibility in (b) also follows from \cite[Proposition~5.1]{BP26}.
The reducibility criterion that turns (a) into (b) is
\cite[Theorem~2.8, Corollary~2.11]{TZ18}. New in this paper are a uniform
proof of (a)--(c) for all $N$ and both exponent lists, which includes the
sizes $N\le10$ stated in \cite{BP26}, and the implication (a)$\Rightarrow$(c)
on unbounded intervals for weights of the form $W=wP$
(Proposition~\ref{prop:bridge}). The latter is needed because the converse
part of \cite[Theorem~4.3]{TZ18} assumes a bounded interval, while the
Laguerre-type weights live on $(0,\infty)$; we therefore prove the required
implication directly for the weights $W=wP$ considered here. Our argument
uses only that $W_s$ is a scalar weight times a polynomial matrix, and
assumes neither bounded support nor determinacy of the moment problem.

\subsection*{Method}
The core of the proof is a direct computation of $\CR(W_s)$. Conjugating a
$W_s$-symmetric matrix $F$ by $T(x)$ and ordering the coordinates as odd
indices followed by even indices, the condition $FW_s=W_sF^*$ becomes the
identity $G(x)D_s(x)=D_s(x)G^*(x)$ between matrix polynomials, where $G^*$
is the coefficientwise adjoint and the lower-left block of $G$ is a constant
matrix $U$. Degree comparisons, together with the full column rank of the
rectangular block of $A$, force $U=0$; the identity then forces $G$ to be
real diagonal, and the path structure of $A$ forces all diagonal entries to
coincide. The repeated exponent $0$ in $\sL_N$ needs one
additional equation from the odd block.

\subsection*{Organization}
Section~\ref{sec:prelim} collects facts on operators of order zero and
proves the implication (a)$\Rightarrow$(c) for weights of the form $W=wP$. Section~\ref{sec:block} introduces the block form, and
Section~\ref{sec:proof} proves Theorem~\ref{thm:main}.

\section{Operators of order zero and \texorpdfstring{$W$}{W}-symmetric matrices}\label{sec:prelim}

Throughout this section $W$ is a weight matrix of size $N$ on an interval
$\Delta$, which may be unbounded, and
$\ip PQ=\int_\Delta P(x)W(x)Q(x)^*\,dx$ for matrix polynomials $P,Q$. If $v$
is a nonzero row vector polynomial, then $v(x)W(x)v(x)^*>0$ for almost every
$x\in\Delta$, since $v$ has finitely many zeros; hence $\int v W v^*>0$. Thus
the block Hankel matrices of the moments of $W$ are positive definite, and
there is a unique sequence $(Q_n)_{n\ge0}$ of monic orthogonal polynomials,
$Q_n(x)=x^nI+\text{(lower terms)}$ and $\ip{Q_n}{Q_m}=0$ for $n\ne m$; see
\cite{Kre49} and \cite[Section~2.1]{BP26}. In particular $Q_0=I$.

\begin{lemma}\label{lem:orderzero}
An operator $D=F_0(x)$ of order zero belongs to $\cD(W)$ if and only if $F_0$
is a constant matrix $C$ such that $CQ_n=Q_nC$ for every $n\ge0$.
\end{lemma}

\begin{proof}
Since $\cD(W)$ does not depend on the sequence of orthogonal polynomials
\cite[Remark~2.3]{BP26}, we use $(Q_n)$. If $D\in\cD(W)$, then
$Q_0\cdot D=F_0(x)=\Lambda_0(D)$ is constant, say $F_0=C$. For $n\ge0$ the
identity $Q_n(x)C=\Lambda_n(D)Q_n(x)$ gives $C=\Lambda_n(D)$ on comparing
the coefficients of $x^n$. Conversely, if $C$ commutes with every $Q_n$, then
$Q_n\cdot C=CQ_n$, so $C\in\cD(W)$ with $\Lambda_n(C)=C$.
\end{proof}

We identify $\cD_0(W)$ with the algebra of constant matrices described in
Lemma~\ref{lem:orderzero}. The next lemma records the facts from
\cite{TZ18} and \cite{KR16} that we use; their proofs do not use the
boundedness of $\Delta$. Part (i) concerns only the values $W(x)$, and the
proof of part (ii) uses only the integral inner product and orthogonality.

\begin{lemma}\label{lem:TZ}
Let $W$ be a weight matrix on an interval $\Delta$, bounded or not.
\begin{enumerate}
\item[(i)] $W$ is irreducible if and only if $\CR(W)=\R I$
\cite[Theorem~2.8, Corollary~2.11]{TZ18}.
\item[(ii)] $\CR(W)\subseteq\cD_0(W)$ \cite[Theorem~4.2]{TZ18},
\cite[Lemma~3.1]{KR16}.
\end{enumerate}
\end{lemma}

\begin{proof}
Part (ii) is the inclusion proved in the first half of the proof of
\cite[Theorem~4.2]{TZ18}; we recall the argument. Let $F\in\CR(W)$, so that
$\ip{PF}{Q}=\ip{P}{QF}$ for all matrix polynomials $P,Q$. Suppose that $F$
commutes with $Q_j$ for $j<n$. Then
$\ip{Q_nF-FQ_n}{Q_j}=\ip{Q_n}{Q_jF}-F\ip{Q_n}{Q_j}=\ip{Q_n}{FQ_j}=0$ for
$j<n$. Since $Q_n$ is monic, $Q_nF-FQ_n$ has degree less than $n$ and is
orthogonal to all polynomials of degree less than $n$; hence it vanishes.
By induction $F$ commutes with every $Q_n$, and $F\in\cD_0(W)$ by
Lemma~\ref{lem:orderzero}.
\end{proof}

On a bounded interval the reverse implication, from a non-scalar element of
$\cD_0(W)$ to a non-scalar element of $\CR(W)$, is
\cite[Theorems~4.2 and~4.3]{TZ18}; its proof uses that a function on a
bounded interval whose moments all vanish is zero. For the weights of this
paper the following replaces that step.

\begin{proposition}\label{prop:bridge}
Let $W=wP$, where $w$ is a scalar weight on an interval $\Delta$, possibly
unbounded, and $P$ is an $N\times N$ matrix polynomial with $P(x)^*=P(x)$
for real $x$, such that $W$ is a weight matrix on $\Delta$. Let
$M_0=\int_\Delta W(x)\,dx$ and $K=M_0^{-1/2}$. If $C\in\cD_0(W)$, then
$KCK^{-1}$ commutes with $KP(x)K$ for every $x\in\R$. Consequently,
$\cD_0(W)=\C I$ if and only if $\CR(W)=\R I$.
\end{proposition}

\begin{proof}
The matrix $M_0$ is positive definite by the remarks at the beginning of
this section. Put $\widetilde P=KPK$ and $\widetilde W=KWK=w\widetilde P$,
so that $\int_\Delta\widetilde W=I$. The polynomials
$R_n=KQ_nK^{-1}$ are monic, and
$\int_\Delta R_n\widetilde WR_m^*=K\ip{Q_n}{Q_m}K$ vanishes for $n\ne m$;
thus $(R_n)$ is the monic orthogonal sequence of $\widetilde W$. The matrix
$C'=KCK^{-1}$ commutes with every $R_n$, hence with every coefficient of
$R_n$. Write $R_n(x)=\sum_{j=0}^nL_{n,j}x^j$ with $L_{n,n}=I$, and let
$\widetilde M_k=\int_\Delta x^k\widetilde W(x)\,dx$. Orthogonality of $R_n$
to $R_0=I$ gives
\[
\widetilde M_n=-\sum_{j=0}^{n-1}L_{n,j}\widetilde M_j\qquad(n\ge1).
\]
Since $\widetilde M_0=I$, induction on $n$ shows that $C'$ commutes with
every $\widetilde M_k$. Let $z$ be an entry of the matrix polynomial
$C'\widetilde P-\widetilde PC'$, say $z(x)=\sum_{j=0}^dc_jx^j$. Then
\begin{equation}\label{eq:moments}
\int_\Delta x^kw(x)z(x)\,dx=0\qquad(k\ge0),
\end{equation}
because the left-hand side is an entry of
$C'\widetilde M_k-\widetilde M_kC'$; all these integrals converge absolutely
because $w$ has finite moments. Multiplying \eqref{eq:moments} for
$k=0,\dots,d$ by $\conj{c_k}$ and adding, we obtain
$\int_\Delta w(x)\lvert z(x)\rvert^2\,dx=0$, where we used that $x$ is
real. As $w>0$ almost everywhere on $\Delta$, the polynomial $z$ vanishes
almost everywhere on the interval $\Delta$ and hence identically. This
proves the first assertion.

If $\CR(W)=\R I$ and $C\in\cD_0(W)$, then $C'$ commutes with the Hermitian
matrices $\widetilde P(x)$, $x\in\R$, and so does $C'^*$. Both
$H_1=(C'+C'^*)/2$ and $H_2=(C'-C'^*)/(2i)$ are Hermitian and commute with
every $\widetilde W(x)$. For $H\in\{H_1,H_2\}$ the matrix $F=K^{-1}HK$
satisfies $F^*=KHK^{-1}$ and therefore
\[
FW=K^{-1}H\widetilde WK^{-1}=K^{-1}\widetilde WHK^{-1}=WF^*,
\]
that is, $F\in\CR(W)=\R I$. Hence $H_1$ and $H_2$ are real multiples of $I$,
and $C=K^{-1}(H_1+iH_2)K\in\C I$. Conversely, if $\cD_0(W)=\C I$, then
Lemma~\ref{lem:TZ}(ii) gives $\CR(W)\subseteq\C I$, and $cW=W\conj c$ with
$W\ne0$ forces $c\in\R$.
\end{proof}

\section{The block form}\label{sec:block}

Let $N\ge2$, put $p=\lceil N/2\rceil$ and $q=\lfloor N/2\rfloor$, and order
the coordinates of $\C^N$ as the odd indices $1,3,\dots,2p-1$ followed by
the even indices $2,4,\dots,2q$. This reordering is a conjugation by a
permutation matrix; it preserves irreducibility and transforms $\CR(W)$ by
the same conjugation, so we may work in the new coordinates. There
\[
A=\begin{pmatrix}0&B\\0&0\end{pmatrix},\qquad
T(x)=\begin{pmatrix}I_p&xB\\0&I_q\end{pmatrix},\qquad
T(x)^{-1}=\begin{pmatrix}I_p&-xB\\0&I_q\end{pmatrix},
\]
where $B$ is the real $p\times q$ matrix whose only nonzero entries are
\begin{equation}\label{eq:B}
B_{r,r}=a_{2r-1}\quad(1\le r\le q),\qquad
B_{r+1,r}=a_{2r}\quad(1\le r\le q,\ r+1\le p).
\end{equation}
Indeed, the entry $a_{2r-1}$ of $A$ sits in row $2r-1$ and column $2r$, and
the entry $a_{2r}$ sits in row $2r+1$ and column $2r$; the condition
$r\le\lfloor(N-1)/2\rfloor$ in \eqref{eq:A} is equivalent to $r+1\le p$.
In particular $B$ is the biadjacency matrix of a path through all $N$
coordinates, with nonzero weights.

\begin{lemma}\label{lem:rank}
The matrix $B$ has rank $q$. Consequently $BY=0$ implies $Y=0$ for every
complex matrix $Y$ with $q$ rows.
\end{lemma}

\begin{proof}
By \eqref{eq:B} the first $q$ rows of $B$ form a lower bidiagonal matrix
with diagonal $a_1,a_3,\dots,a_{2q-1}$, whose determinant is nonzero.
\end{proof}

In the new coordinates $D_s=\diag(D_o,D_e)$ with
$D_o=\diag(x^{o_1},\dots,x^{o_p})$ and $D_e=\diag(x^{e_1},\dots,x^{e_q})$,
where $o_r=s_{2r-1}$ and $e_t=s_{2t}$. Explicitly,
\begin{equation}\label{eq:exponents}
\begin{aligned}
s=\sE_N&\colon\quad o_r=2r-2,\qquad e_t=2t-1;\\
s=\sL_N&\colon\quad o_1=0,\quad o_r=2r-3\ (r\ge2),\qquad e_t=2t-2.
\end{aligned}
\end{equation}

\begin{lemma}\label{lem:exponents}
For $s\in\{\sL_N,\sE_N\}$:
\begin{enumerate}
\item[(i)] $o_1,\dots,o_p$ are pairwise distinct, and so are
$e_1,\dots,e_q$;
\item[(ii)] for $1\le r\le p$ and $1\le t\le q$ we have
$o_r-e_t\in\{0,1,2\}$ if and only if either $r=t+1$, in which case
$o_r-e_t=1$, or $s=\sL_N$ and $r=t=1$, in which case $o_r-e_t=0$.
\end{enumerate}
\end{lemma}

\begin{proof}
Part (i) is clear from \eqref{eq:exponents}. In both cases
$o_r-e_t=2(r-t)-1$ unless $s=\sL_N$ and $r=1$; this odd number lies in
$\{0,1,2\}$ exactly when $r=t+1$. If $s=\sL_N$ and $r=1$, then
$o_1-e_t=2-2t$, which lies in $\{0,1,2\}$ exactly when $t=1$.
\end{proof}

Let $F\in\Mat_N(\C)$, written in blocks of sizes $p$ and $q$ as
\[
F=\begin{pmatrix}R&Y\\U&S\end{pmatrix}.
\]
Multiplying the triangular block matrices gives
\begin{equation}\label{eq:G}
G(x)=T(x)^{-1}FT(x)=
\begin{pmatrix}
R-xBU & Y+x(RB-BS)-x^2BUB\\
U & S+xUB
\end{pmatrix}.
\end{equation}
For a matrix polynomial $Z(x)=\sum_jZ_jx^j$ we write
$Z^*(x)=\sum_jZ_j^*x^j$, so that $Z^*(x)=Z(x)^*$ for real $x$.

\begin{lemma}\label{lem:polyidentity}
If $F\in\CR(W_s)$, then
\begin{equation}\label{eq:GD}
G(x)D_s(x)=D_s(x)G^*(x)
\end{equation}
as an identity of matrix polynomials.
\end{lemma}

\begin{proof}
Since $w>0$ almost everywhere on $\Delta$, the relation $FW_s=W_sF^*$ gives
$FP_s(x)=P_s(x)F^*$ for almost every $x\in\Delta$. Both sides are matrix
polynomials, and a polynomial vanishing on a set of positive measure is zero,
so this is an identity of matrix polynomials. Since $T$ has real
coefficients, $T^*=T^{\mathsf T}$ and $(T^{-1}FT)^*=T^*F^*(T^*)^{-1}$;
multiplying the identity on the left by $T^{-1}$ and on the right by
$(T^*)^{-1}$ gives \eqref{eq:GD}.
\end{proof}

In blocks, \eqref{eq:GD} consists of the three identities
\begin{gather}
(R-xBU)D_o=D_o(R-xBU)^*,\label{eq:oo}\\
\bigl(Y+x(RB-BS)-x^2BUB\bigr)D_e=D_oU^*,\label{eq:oe}\\
(S+xUB)D_e=D_e(S+xUB)^*,\label{eq:ee}
\end{gather}
the fourth block being the adjoint of \eqref{eq:oe}.

\section{Proof of Theorem~\ref{thm:main}}\label{sec:proof}

\begin{proposition}\label{prop:CR}
Under the hypotheses of Theorem~\ref{thm:main}, $\CR(W_s)=\R I$.
\end{proposition}

We first treat $N=1$. Then $W_s=w$ is scalar, and $cw=w\conj c$ forces
$c\in\R$, so $\CR(W_s)=\R I$. From now on let $N\ge2$ and
$F\in\CR(W_s)$, with the notation of Section~\ref{sec:block}. We show that
the lower-left block of $G$ vanishes.

\begin{lemma}\label{lem:U}
$U=0$.
\end{lemma}

\begin{proof}
Write $G_{12}=Y+x(RB-BS)-x^2BUB$ for the upper-right block of $G$, a
matrix polynomial of degree at most two. The $(r,t)$ entry of
\eqref{eq:oe} reads
\begin{equation}\label{eq:entry}
(G_{12})_{r,t}(x)\,x^{e_t}=\conj{U_{t,r}}\,x^{o_r}.
\end{equation}
If $U_{t,r}\ne0$, the right-hand side is a nonzero monomial divisible by
$x^{e_t}$, so $o_r\ge e_t$ and $(G_{12})_{r,t}=\conj{U_{t,r}}x^{o_r-e_t}$;
as $G_{12}$ has degree at most two, $o_r-e_t\in\{0,1,2\}$. By
Lemma~\ref{lem:exponents}(ii) the only entries of $U$ that can be nonzero
are
\begin{equation}\label{eq:Uentries}
u_t=U_{t,t+1}\quad(1\le t\le q,\ t+1\le p),\qquad
v=U_{1,1}\quad\text{(only if $s=\sL_N$)},
\end{equation}
and for each of them $(G_{12})_{r,t}$ has degree at most one. All other
entries of $G_{12}$ vanish by \eqref{eq:entry}. Hence the coefficient of
$x^2$ in $G_{12}$ vanishes, that is, $BUB=0$, and Lemma~\ref{lem:rank}
gives
\begin{equation}\label{eq:UB}
UB=0.
\end{equation}

Let $1\le t\le q$ with $t+1\le p$. By \eqref{eq:B} and
\eqref{eq:Uentries},
\begin{equation}\label{eq:UBdiag}
(UB)_{t,t}=U_{t,t}a_{2t-1}+u_ta_{2t},
\end{equation}
where $U_{t,t}=0$ unless $s=\sL_N$ and $t=1$.

\emph{Case $s=\sE_N$.} Here $U_{t,t}=0$ for all $t$, so \eqref{eq:UB} and
\eqref{eq:UBdiag} give $u_ta_{2t}=0$, hence $u_t=0$, for every $t$ in
\eqref{eq:Uentries}. Thus $U=0$.

\emph{Case $s=\sL_N$, $N=2$.} Here $p=q=1$, the only possible entry is $v$,
and \eqref{eq:UB} reads $va_1=0$. Thus $U=0$.

\emph{Case $s=\sL_N$, $N\ge3$.} Now $p\ge2$, $o_1=0$ and $o_2=1$. The
$(2,1)$ entry of \eqref{eq:oo} is
\[
(R-xBU)_{2,1}(x)=x\,\conj{(R-xBU)_{1,2}}(x)
=x\,\conj{R_{1,2}}-x^2\,\conj{(BU)_{1,2}}.
\]
The left-hand side has degree at most one, so $(BU)_{1,2}=0$. The first row
of $B$ has the single nonzero entry $B_{1,1}=a_1$, and $U_{1,2}=u_1$, so
$a_1u_1=0$ and $u_1=0$. Now \eqref{eq:UB} and \eqref{eq:UBdiag} with $t=1$
give $va_1=0$, so $v=0$. For $t\ge2$ we have $U_{t,t}=0$, and as in the case
$s=\sE_N$ we get $u_t=0$. Thus $U=0$.
\end{proof}

\begin{proof}[Proof of Proposition~\ref{prop:CR}]
By Lemma~\ref{lem:U}, $G_{12}=Y+x(RB-BS)$ and \eqref{eq:oe} becomes
$G_{12}D_e=0$. Since the diagonal entries of $D_e$ are nonzero monomials,
$G_{12}=0$, that is,
\begin{equation}\label{eq:RBBS}
Y=0,\qquad RB=BS.
\end{equation}
With $U=0$, identity \eqref{eq:oo} reads $R_{r,t}x^{o_t}=\conj{R_{t,r}}x^{o_r}$
for all $r,t$. For $r\ne t$ the exponents differ by
Lemma~\ref{lem:exponents}(i), so $R_{r,t}=0$; for $r=t$ we get
$R_{r,r}\in\R$. Thus $R=\diag(\rho_1,\dots,\rho_p)$ is real, and in the same
way \eqref{eq:ee} shows that $S=\diag(\sigma_1,\dots,\sigma_q)$ is real.
Comparing the entries of $RB=BS$ at the nonzero entries \eqref{eq:B} of
$B$ gives
\[
\rho_t=\sigma_t\quad(1\le t\le q),\qquad
\rho_{t+1}=\sigma_t\quad(1\le t\le q,\ t+1\le p).
\]
Starting from $\rho_1=\sigma_1$ and using the two relations alternately we
obtain $\rho_1=\sigma_1=\rho_2=\sigma_2=\cdots$, which reaches $\rho_p$ and
$\sigma_q$. Hence $F=G=\rho_1I$ with $\rho_1\in\R$. Conversely, every real
multiple of $I$ lies in $\CR(W_s)$.
\end{proof}

\begin{proof}[Proof of Theorem~\ref{thm:main}]
Part (a) is Proposition~\ref{prop:CR}, and part (b) follows from (a) and
Lemma~\ref{lem:TZ}(i). For part (c), $W_s=wP_s$, where $w$ is a scalar weight
on $\Delta$ because $s_1=0$, and $P_s=TD_sT^*$ is a matrix polynomial with
$P_s(x)^*=P_s(x)$ for real $x$. Proposition~\ref{prop:bridge} and part (a)
give $\cD_0(W_s)=\C I$.
\end{proof}

\begin{proof}[Proof of Corollary~\ref{cor:conjecture}]
Put $w=w_1$. In the first case $w_j=w\,x^{s_j}$ with $s=\sL_N$, in the
second with $s=\sE_N$, and $w$ is $s$-admissible because each $w_j$ is a
scalar weight on the common interval. The weight \eqref{eq:BPweight} is
$W_s$, and Theorem~\ref{thm:main} applies.
\end{proof}

\begin{remark}
The proof of Proposition~\ref{prop:CR} compares coefficients and degrees in
identities of matrix polynomials, so it does not require that $0$ lie in
$\Delta$. The scalar weight $w$ enters Theorem~\ref{thm:main} only through
Lemma~\ref{lem:polyidentity} and Proposition~\ref{prop:bridge}, and neither
uses determinacy of the moment problem.
\end{remark}

\section*{Use of AI-assisted tools}
AI-assisted tools were used in the research and preparation of this manuscript.

\begin{thebibliography}{99}

\bibitem{BP26}
I.~Bono Parisi, \emph{Explicit construction of matrix-valued orthogonal
polynomials of arbitrary size}, J. Spectr. Theory \textbf{16} (2026), no.~1,
299--328.

\bibitem{BPP}
I.~Bono Parisi and I.~Pacharoni, \emph{Singular solutions of the matrix
Bochner problem: the $N$-dimensional cases}, preprint (2024),
arXiv:2411.00798v1.

\bibitem{GT07}
F.~A. Gr\"unbaum and J.~Tirao, \emph{The algebra of differential operators
associated to a weight matrix}, Integral Equations Operator Theory
\textbf{58} (2007), no.~4, 449--475.

\bibitem{KR16}
E.~Koelink and P.~Rom\'an, \emph{Orthogonal vs.\ non-orthogonal
reducibility of matrix-valued measures}, SIGMA Symmetry Integrability Geom.
Methods Appl. \textbf{12} (2016), Paper No.~008, 9~pp.

\bibitem{Kre49}
M.~G. Kre\u{\i}n, \emph{Infinite $J$-matrices and a matrix-moment problem},
Doklady Akad. Nauk SSSR (N.S.) \textbf{69} (1949), 125--128 (Russian).

\bibitem{TZ18}
J.~Tirao and I.~Zurri\'an, \emph{Reducibility of matrix weights}, Ramanujan
J. \textbf{45} (2018), no.~2, 349--374; numbered results are cited from
arXiv:1501.04059v5.

\end{thebibliography}

\end{document}
